\documentclass{article}
\usepackage{amsmath,amssymb,amsthm}
\usepackage{algorithm}
\usepackage{algorithmic}
\usepackage{hyperref}
\usepackage{geometry}
\geometry{margin=1in}
\newtheorem{theorem}{Theorem}
\newtheorem{lemma}{Lemma}
\newcommand{\E}{\mathbb{E}}
\newcommand{\R}{\mathbb{R}}
\newcommand{\norm}[1]{\left\lVert #1\right\rVert}
\begin{document}

%====================================================================
\section*{Adaptive Gradient Descent (AdaGrad) with Cumulative Gradient Sum}
%====================================================================

%--------------------------------------------------------------------
\section*{Mathematical Formulation}
%--------------------------------------------------------------------
Let $f:\mathbb{R}^d\rightarrow\mathbb{R}$ be a convex and $L$‑smooth objective.
Denote by $g_k:=\nabla f(w_k)$ the (full) gradient at iteration $k$.
Define the accumulated squared‑gradient norm
\[
G_k \;:=\;\sqrt{\sum_{i=1}^{k}\|g_i\|^2},\qquad G_0:=0 .
\]
Given a base stepsize $\alpha>0$, the adaptive stepsize at iteration $k$ is
\[
\eta_k \;:=\;\frac{\alpha}{G_k}\quad\text{(with the convention }\eta_1:=\alpha\text{)} .
\]
The update rule (without projection) reads
\begin{equation}
\boxed{%
w_{k+1}= w_k-\eta_k\,g_k
      = w_k-\frac{\alpha}{G_k}\,g_k ,\qquad k\ge 1.}
\label{eq:update}
\end{equation}
If a feasible set $\mathcal{W}\subset\mathbb{R}^d$ is prescribed,
the iterate is projected back:
\[
w_{k+1}= \Pi_{\mathcal{W}}\!\bigl(w_k-\eta_k g_k\bigr).
\]

%--------------------------------------------------------------------
\section*{Pseudocode}
%--------------------------------------------------------------------
\begin{algorithm}[H]
\caption{AdaGrad with cumulative gradient sum}
\label{alg:adagrad}
\begin{algorithmic}[1]
\STATE \textbf{Input:} initial point $w_0\in\mathcal{W}$, base stepsize $\alpha>0$
\STATE $G_0\gets 0$
\FOR{$k=1,2,\dots,T$}
    \STATE Compute gradient $g_k \gets \nabla f(w_{k-1})$
    \STATE $G_k \gets \sqrt{G_{k-1}^2+\|g_k\|^{2}}$
    \STATE $\eta_k \gets \alpha / G_k$   \COMMENT{adaptive stepsize}
    \STATE $w_k \gets \Pi_{\mathcal{W}}\bigl(w_{k-1}-\eta_k g_k\bigr)$
\ENDFOR
\STATE \textbf{Return} $w_T$ (or the average $\bar w_T=\frac1T\sum_{k=1}^{T}w_k$)
\end{algorithmic}
\end{algorithm}

%--------------------------------------------------------------------
\section*{Dry Run Example}
%--------------------------------------------------------------------
Consider the one–dimensional quadratic
\[
f(w)=(w-3)^2,\qquad \nabla f(w)=2(w-3).
\]
Set $w_0=0$, base stepsize $\alpha=0.1$, and run three iterations.

\begin{center}
\begin{tabular}{c|c|c|c|c}
$k$ & $w_{k-1}$ & $g_k=2(w_{k-1}-3)$ & $G_k$ & $w_k$ \\ \hline
1 &
$0$ &
$-6$ &
$\sqrt{(-6)^2}=6$ &
$0-\frac{0.1}{6}(-6)=0.1$ \\[4pt]
2 &
$0.1$ &
$-5.8$ &
$\sqrt{6^2+5.8^2}= \sqrt{36+33.64}=8.345$ &
$0.1-\frac{0.1}{8.345}(-5.8)=0.1695$ \\[4pt]
3 &
$0.1695$ &
$-5.661$ &
$\sqrt{8.345^2+5.661^2}= \sqrt{69.64+32.07}=10.085$ &
$0.1695-\frac{0.1}{10.085}(-5.661)=0.2225$
\end{tabular}
\end{center}

Corresponding objective values:
\[
f(w_1)= (0.1-3)^2=8.41,\quad
f(w_2)= (0.1695-3)^2\approx8.04,\quad
f(w_3)= (0.2225-3)^2\approx7.70.
\]
The algorithm makes progress toward the minimiser $w^\star=3$ while the stepsize shrinks as $G_k$ grows.

%====================================================================
\section*{Convergence Theory Development}
%====================================================================

%--------------------------------------------------------------------
\subsection*{Assumptions}
%--------------------------------------------------------------------
\begin{enumerate}
\item \textbf{Convexity.} $f$ is convex on a closed convex set $\mathcal{W}\subset\R^d$:
\[
f(u)\ge f(v)+\langle\nabla f(v),u-v\rangle,\qquad\forall u,v\in\mathcal{W}.
\tag{A1}
\]
\item \textbf{Smoothness.} $f$ is $L$‑smooth:
\[
\|\nabla f(u)-\nabla f(v)\|\le L\|u-v\|,\qquad\forall u,v\in\mathcal{W}.
\tag{A2}
\]
\item \textbf{Bounded domain.} There exists $D>0$ such that
\[
\|w-w^\star\|\le D,\qquad\forall w\in\mathcal{W},
\]
where $w^\star\in\arg\min_{w\in\mathcal{W}}f(w)$.
\item \textbf{Bounded gradients.} $\|g_k\|=\|\nabla f(w_{k-1})\|\le G$ for all $k$.
\end{enumerate}

%--------------------------------------------------------------------
\subsection*{Main Convergence Theorem}
%--------------------------------------------------------------------
\begin{theorem}[AdaGrad convergence for convex smooth objectives]\label{thm:main}
Under (A1)–(A4) and with the iterates generated by \eqref{eq:update},
for any $T\ge 1$ the averaged iterate
\[
\bar w_T:=\frac1T\sum_{k=1}^{T}w_k
\]
satisfies
\[
\boxed{%
f(\bar w_T)-f(w^\star)\;\le\;
\frac{D^2}{2\alpha\,\sqrt{T}}+\frac{\alpha G^2}{2\sqrt{T}} }.
\tag{1}
\]
Choosing $\alpha=D/G$ yields the optimal $O\!\left(\frac{DG}{\sqrt{T}}\right)$ rate.
\end{theorem}

%--------------------------------------------------------------------
\subsection*{Theoretical Properties}
%--------------------------------------------------------------------
\begin{itemize}
\item \textbf{Stability.} The adaptive stepsize $\eta_k=\alpha/G_k$ never exceeds $\alpha$ and monotonically decreases, guaranteeing bounded updates.
\item \textbf{Hyperparameter sensitivity.} The bound \eqref{eq:update} depends on $\alpha$ only via the product $\alpha G$, thus the optimal choice $\alpha=D/G$ is independent of the unknown smoothness constant $L$.
\item \textbf{Comparison to SGD.} Standard (fixed‑stepsize) SGD with $\eta=\Theta(1/\sqrt{T})$ achieves the same $O(1/\sqrt{T})$ rate but requires prior knowledge of $T$. AdaGrad adapts automatically to the observed gradient magnitudes.
\item \textbf{Effect of smoothness.} When $L$ is small, the gradients tend to decay faster, making $G_k$ grow slower and the effective stepsize larger, which can improve practical performance while preserving the worst‑case bound.
\end{itemize}

%====================================================================
\section*{Complete Convergence Proof}
%====================================================================

%--------------------------------------------------------------------
\subsection*{Preliminaries (Definitions and Lemmas)}
%--------------------------------------------------------------------
\begin{lemma}[Three‑point identity]\label{lem:three-point}
For any $u,v\in\mathcal{W}$ and any $\eta>0$,
\[
\|u-\eta g -v\|^{2}= \|u-v\|^{2}
-2\eta\langle g, u-v\rangle
+\eta^{2}\|g\|^{2}.
\tag{2}
\]
\end{lemma}
\begin{proof}
Expand the square and collect terms; the identity holds for any vectors.
\end{proof}

\begin{lemma}[Bounding the adaptive stepsizes]\label{lem:stepsize}
Let $G_k$ be defined as above and $\eta_k=\alpha/G_k$.
Then for every $k\ge 1$,
\[
\eta_k\|g_k\|^{2}
   =\alpha\frac{\|g_k\|^{2}}{G_k}
   \le \alpha\bigl(G_k-G_{k-1}\bigr).
\tag{3}
\]
\end{lemma}
\begin{proof}
Since $G_k^{2}=G_{k-1}^{2}+\|g_k\|^{2}$,
\[
G_k-G_{k-1}
    =\frac{G_k^{2}-G_{k-1}^{2}}{G_k+G_{k-1}}
    =\frac{\|g_k\|^{2}}{G_k+G_{k-1}}
    \ge\frac{\|g_k\|^{2}}{2G_k}.
\]
Multiplying both sides by $2\alpha$ yields \eqref{3}.
\end{proof}

%--------------------------------------------------------------------
\subsection*{Main Proof (Step‑by‑Step Derivation)}
%--------------------------------------------------------------------
We prove Theorem~\ref{thm:main}. Throughout the proof we fix an arbitrary optimal point $w^\star\in\mathcal{W}$.

\noindent\textbf{[Step 1]} Apply Lemma~\ref{lem:three-point} with $u=w_{k-1}$, $v=w^\star$, $g=g_{k}$ and $\eta=\eta_k$:
\[
\|w_k-w^\star\|^{2}
   =\|w_{k-1}-w^\star\|^{2}
   -2\eta_k\langle g_k,w_{k-1}-w^\star\rangle
   +\eta_k^{2}\|g_k\|^{2}.
\tag{4}
\]

\noindent\textbf{[Step 2]} Rearrange \eqref{4} to isolate the inner product term:
\[
2\eta_k\langle g_k,w_{k-1}-w^\star\rangle
   =\|w_{k-1}-w^\star\|^{2}
    -\|w_k-w^\star\|^{2}
    +\eta_k^{2}\|g_k\|^{2}.
\tag{5}
\]

\noindent\textbf{[Step 3]} Sum \eqref{5} over $k=1,\dots,T$:
\[
2\sum_{k=1}^{T}\eta_k\langle g_k,w_{k-1}-w^\star\rangle
   =\|w_{0}-w^\star\|^{2}
    -\|w_T-w^\star\|^{2}
    +\sum_{k=1}^{T}\eta_k^{2}\|g_k\|^{2}.
\tag{6}
\]

\noindent\textbf{[Step 4]} Use the bound $\|w_0-w^\star\|\le D$ and discard the non‑negative term $\|w_T-w^\star\|^{2}$:
\[
2\sum_{k=1}^{T}\eta_k\langle g_k,w_{k-1}-w^\star\rangle
   \le D^{2}
    +\sum_{k=1}^{T}\eta_k^{2}\|g_k\|^{2}.
\tag{7}
\]

\noindent\textbf{[Step 5]} By convexity (A1),
\[
f(w_{k-1})-f(w^\star)
   \le \langle g_k, w_{k-1}-w^\star\rangle .
\tag{8}
\]
Multiplying \eqref{8} by $\eta_k$ and summing yields
\[
\sum_{k=1}^{T}\eta_k\bigl[f(w_{k-1})-f(w^\star)\bigr]
   \le \sum_{k=1}^{T}\eta_k\langle g_k,w_{k-1}-w^\star\rangle .
\tag{9}
\]

\noindent\textbf{[Step 6]} Combine \eqref{7} and \eqref{9}:
\[
2\sum_{k=1}^{T}\eta_k\bigl[f(w_{k-1})-f(w^\star)\bigr]
   \le D^{2}
    +\sum_{k=1}^{T}\eta_k^{2}\|g_k\|^{2}.
\tag{10}
\]

\noindent\textbf{[Step 7]} Apply Lemma~\ref{lem:stepsize} to bound the second sum:
\[
\eta_k^{2}\|g_k\|^{2}
   =\eta_k\bigl(\eta_k\|g_k\|^{2}\bigr)
   \le\eta_k\alpha\bigl(G_k-G_{k-1}\bigr)
   =\alpha\frac{\alpha}{G_k}\bigl(G_k-G_{k-1}\bigr)
   =\alpha^{2}\frac{G_k-G_{k-1}}{G_k}.
\tag{11}
\]
Since $G_k\ge G_{k-1}$, we have $\frac{G_k-G_{k-1}}{G_k}\le \frac{G_k-G_{k-1}}{G_{k-1}+G_k\le 1$, and a simpler bound
\[
\eta_k^{2}\|g_k\|^{2}\le \alpha\bigl(G_k-G_{k-1}\bigr).
\tag{12}
\]

\noindent\textbf{[Step 8]} Summing \eqref{12} telescopes:
\[
\sum_{k=1}^{T}\eta_k^{2}\|g_k\|^{2}
   \le \alpha\sum_{k=1}^{T}\bigl(G_k-G_{k-1}\bigr)
   = \alpha\,(G_T-G_0)=\alpha\,G_T .
\tag{13}
\]

\noindent\textbf{[Step 9]} Insert \eqref{13} into \eqref{10}:
\[
2\sum_{k=1}^{T}\eta_k\bigl[f(w_{k-1})-f(w^\star)\bigr]
   \le D^{2}+\alpha G_T .
\tag{14}
\]

\noindent\textbf{[Step 10]} Because $\eta_k=\alpha/G_k\ge \alpha/G_T$ for all $k$, we have
\[
\eta_k\ge \frac{\alpha}{G_T}\quad\Longrightarrow\quad
\sum_{k=1}^{T}\eta_k
   \ge T\frac{\alpha}{G_T}.
\tag{15}
\]

\noindent\textbf{[Step 11]} Divide \eqref{14} by $\sum_{k=1}^{T}\eta_k$ and use \eqref{15}:
\[
\frac{\displaystyle\sum_{k=1}^{T}\eta_k\bigl[f(w_{k-1})-f(w^\star)\bigr]}
     {\displaystyle\sum_{k=1}^{T}\eta_k}
   \;\le\;
   \frac{D^{2}+\alpha G_T}{2\,\sum_{k=1}^{T}\eta_k}
   \;\le\;
   \frac{D^{2}+\alpha G_T}{2\,T\alpha/G_T}
   =\frac{D^{2}G_T}{2\alpha T}+\frac{G_T^{2}}{2T}.
\tag{16}
\]

\noindent\textbf{[Step 12]} By the bounded‑gradient assumption $\|g_k\|\le G$, we have $G_T\le G\sqrt{T}$. Substituting this bound into \eqref{16} yields
\[
\frac{1}{T}\sum_{k=1}^{T}\bigl[f(w_{k-1})-f(w^\star)\bigr]
   \le \frac{D^{2}G}{2\alpha\sqrt{T}}+\frac{\alpha G^{2}}{2\sqrt{T}}.
\tag{17}
\]

\noindent\textbf{[Step 13]} By convexity, the average of the function values dominates the function value at the average iterate:
\[
f(\bar w_T)-f(w^\star)
   \le \frac{1}{T}\sum_{k=1}^{T}\bigl[f(w_{k-1})-f(w^\star)\bigr].
\tag{18}
\]

\noindent\textbf{[Step 14]} Combining \eqref{18} and \eqref{17} gives exactly the bound claimed in \eqref{1}:
\[
f(\bar w_T)-f(w^\star)
   \le \frac{D^{2}}{2\alpha\sqrt{T}}+\frac{\alpha G^{2}}{2\sqrt{T}}.
\tag{19}
\]
\hfill $\square$

%--------------------------------------------------------------------
\subsection*{Rate Derivation (With Constants)}
%--------------------------------------------------------------------
The right‑hand side of \eqref{19} is minimized by choosing
\[
\alpha^\star = \frac{D}{G},
\]
which yields
\[
f(\bar w_T)-f(w^\star)
   \le \frac{DG}{\sqrt{T}} .
\tag{20}
\]
Thus AdaGrad attains the optimal $O(1/\sqrt{T})$ convergence rate for convex smooth problems without any knowledge of $T$.

%--------------------------------------------------------------------
\subsection*{Special Cases}
%--------------------------------------------------------------------
\begin{itemize}
\item \textbf{Strongly convex $f$.} If $f$ is $\mu$‑strongly convex, a standard modification (averaging the iterates with weights $\eta_k$) improves the rate to $O\!\bigl(\frac{DG}{\mu T}\bigr)$.
\item \textbf{Sparse gradients.} When many coordinates of $g_k$ are zero, the per‑coordinate version of AdaGrad yields $G_T=\sqrt{\sum_{j=1}^{d}\sum_{k=1}^{T}g_{k,j}^{2}}$, often growing like $\sqrt{\log T}$, leading to a tighter $O\!\bigl(\frac{DG\sqrt{\log T}}{\sqrt{T}}\bigr)$ bound.
\item \textbf{Deterministic gradients (full‑batch).} The analysis above does not rely on stochasticity; therefore the same bound holds for deterministic gradient descent with the adaptive stepsize.
\end{itemize}

%====================================================================
\section*{Discussion}
%====================================================================
The proof shows that the adaptive denominator $G_k$ automatically balances the magnitude of the updates. No explicit decay schedule is required; the algorithm self‑regulates the learning rate. The derived $O(1/\sqrt{T})$ bound matches the optimal rate for convex Lipschitz functions and improves over fixed‑step SGD when the observed gradients are sparse or decay quickly. Moreover, the analysis hinges only on convexity, smoothness, bounded domain, and bounded gradients—exactly the standard AdaGrad assumptions.

\end{document}